%% kapa: exam=1 round=23 year=2012 period=1 form=A subject=civil_law subject_ko=민법 seq=1 number=1
%% source: Q-Net 감정평가사 기출문제 (cst00309) artlSeq=5052578 file=raw/1008/5052578/2012년 제23회 감정평가사 제1차시험 기출문제(1교시 A형).pdf
%% conversion: kordoc 4.24.51 -> tools/gen1.py (auto); review=proofread-sonnet; flags=none
\begin{question}{1}
물건에 관한 설명으로 옳은 것은? (다툼이 있으면 판례에 의함)
\begin{choices}
  \item 완공된 신축 건물이라도 보존등기 전에는 소유권이 인정되지 않는다.
  \item 1필의 토지 일부에는 지상권을 설정할 수 없다.
  \item 등기된 입목의 소유자는 입목을 토지와 분리하여 양도할 수 있고 이를 저당권의 목적으로 할 수 있다.
  \item 권한 없이 타인의 토지에 농작물을 심은 경우, 수확기에 이른 농작물의 소유권은 토지 소유자에게 귀속된다.
  \item 주유소의 주유기는 주유소 건물에 부합되어 종물로 볼 수 없다.
\end{choices}
\end{question}
\answer{3}
